The residual baseline is not a fixed prefactor arbitrarily chosen and then fed into the model.

It is treated as an outcome derived from the factual 2026 macro ledger:

U.S. federal deficit (CBO) ≈ $1.90 T  
Global structural / off-balance-sheet / SFA claims ≈ $1.75 T  

→ Combined factual annual residual base of $3.65 T.

This $3.65 T figure is therefore an empirical outcome of the current claim-issuance structure (after applying the low true-repayment fraction \(\rho_t\)), not an exogenous tuning parameter.

Consequence for the Timeline Claims

Because the residual baseline is itself an outcome of the observed deficit and structural claim data, the subsequent results — median breach at 2055 and full saturation by 2059 — are also outcomes of that same factual grounding once the locked dynamics (regime \(\psi\), reciprocal mirror, low \(\rho_t\), etc.) are applied.

The $25–45 T liquidation band remains a model-generated magnitude that follows from the scale of the residual stock at the crossing point and the endogenous cascade operators.

The residual baseline is repositioned as a derived factual quantity rather than a free prefactor.
